\section{Does giving up exactness pay?}
\label{sec:relaxation}

% Generated by tools/gen_tables.py; do not edit by hand.
\newcommand{\varexact}{4.13}\newcommand{\varfast}{0.36}\newcommand{\encexact}{14.5}\newcommand{\encfast}{10.9}


The fast mode is free to reorder operations, fold symmetric filter taps (halving the
multiplications), use fused multiply-add and compute in double precision where that is cheaper.
Table~\ref{tab:relax} compares it with the exact mode stage by stage.

% Generated by tools/gen_tables.py; do not edit by hand.
\begin{table}[ht]
  \centering\small
  \begin{tabular}{@{}llrrr@{}}
    \toprule
    & stage & exact (ms) & fast (ms) & fast / exact\\
    \midrule
    decoder & entropy & 15 & 15 & 1.00\\
     & dequantize & 10 & 10 & 1.00\\
     & transform & 80 & 81 & 1.01\\
     & output & 8 & 8 & 0.94\\
     & total & 119 & 120 & 1.00\\
    \midrule
    encoder & normalize & 12 & 11 & 0.98\\
     & transform & 106 & 102 & 0.96\\
     & variance & 63 & 9 & 0.14\\
     & quantize & 17 & 17 & 1.00\\
     & total & 242 & 186 & 0.77\\
    \bottomrule
  \end{tabular}
  \caption{Exact against fast arithmetic, AVX2, summed over the 40 images at 0.75 bits per pixel (one core).}
  \label{tab:relax}
\end{table}


\paragraph*{Decoding: no.} With AVX2 the two modes take the same time to the percent. The
transform is limited by data movement (Section~\ref{sec:optimizations}), so doing half the
multiplications changes nothing; dequantization and pixel conversion are memory-bound as well.
Only the scalar backend, which is compute-bound, gains (11\% over the 40 images,
Table~\ref{tab:decode}). Since exactness costs nothing in the configurations that matter
(AVX2 natively, SIMD in WebAssembly), the decoder defaults to exact mode.

\paragraph*{Encoding: yes, in one place.} The encoder is 23\% faster in fast mode over the set,
and almost all of the difference is the subband variances. NBIS sums them in float, one
coefficient after the other; reproducing its rounding forces the same sequential sum, a chain
of dependent additions that no vector unit can shorten. The fast mode sums in double with
independent accumulators. On images whose variance is estimated over whole subbands (7 of the
40, after the version 3.1 rule), the difference is large: \texttt{sample\_19} spends
\varexact~ms on variances in exact mode and \varfast~ms in fast mode, and encodes in \encfast~ms
instead of \encexact~ms. The fast encoder's files meet every encoder criterion of Annex AA
against NBIS.

\paragraph*{What else could a relaxed mode do?} Lifting would cut the arithmetic of the 9/7
transform further, but arithmetic is not the bottleneck. Storing the intermediate of each split
in 16-bit floats would halve the data movement that is, at the cost of precision; that and
integer arithmetic remain untested.
